\section{Περιορισμοί και Υποθέσεις}

Η εργασία επικεντρώνεται στις εκδόσεις \ENG{SCORM 1.2} και \ENG{SCORM 2004}, καθώς αποτελούν τα πιο διαδεδομένα πρότυπα
σε υφιστάμενες πλατφόρμες \ENG{E-Learning}. Άλλα πρότυπα όπως \ENG{xAPI}, \ENG{cmi5} και \ENG{AICC} εξετάζονται
μόνο θεωρητικά και δεν αποτελούν μέρος της υλοποίησης.

Το σύστημα που παρουσιάζεται δεν αποτελεί πλήρες \ENG{Learning Management System}. Τόσο το \ENG{\texttt{SCORM engine}} 
όσο και ο \ENG{\texttt{player}} λειτουργούν ως εξειδικευμένη υποδομή εκτέλεσης και παρακολούθησης μαθημάτων, ενώ λειτουργίες
όπως διαχείριση εκπαιδευτικών προγραμμάτων, κοινωνικές δυνατότητες ή οικονομικές συναλλαγές παραμένουν εκτός του πεδίου της εργασίας.
Το ενδεικτικό \ENG{LMS} χρησιμοποιείται αποκλειστικά ως παράδειγμα \ENG{integration}, επίδειξη δηλαδή της ενσωμάτωσης σε άλλα συστήματα.

Υποτίθεται επίσης ότι τα πακέτα μαθημάτων που εισάγονται είναι έγκυρα \ENG{SCORM} αρχεία. Η υλοποίηση ελέγχει βασικές συνθήκες,
όπως την ύπαρξη του \ENG{\texttt{imsmanifest.xml}} και την αναγνώριση του σημείου εκκίνησης, χωρίς να επιχειρεί πλήρη σημασιολογική
επικύρωση όλων των πιθανών παραλλαγών περιεχομένου.

Τέλος, η υποδομή που παρουσιάζεται βασίζεται σε \ENG{Docker Compose} και είναι κατάλληλη για ανάπτυξη και δοκιμή του συστήματος.
Ζητήματα όπως εκτεταμένη ανθεκτικότητα υποδομής, \ENG{high availability}, κατανεμημένες αναπτύξεις ή πλήρη συστήματα \ENG{learning analytics}
δεν αποτελούν αντικείμενο της παρούσας εργασίας.